\section{Fields, indices, and growth along shortcuts}
\label{grw:section}

We continue the reduction of Theorem~\ref{cls:chain-rigidity}, starting
with the configurations supplied by Proposition~\ref{cls:reduction}.
Their distinct classical endpoint types have arbitrarily large rank,
their multiplicities agree, and their killed kernels are soluble.
We now compare the numerical restrictions imposed by a two-edge shortcut
with the growth obtained by traversing several successive comparisons.
The three quantities are field exponents, natural dimensions, and inner
order quotients. We first establish the maximal-step estimates needed
for these comparisons. They will exclude wholly geometric paths and
characteristic crossings, leaving common-field paths with quadratic
dimension growth and controlled inner order.

Write $z_0<\cdots<z_\ell$ for the retained original vertices, in their
original order, and abbreviate $F_i=F_{z_i}$, $N_i=N_{z_i}$,
$T_i=T_{z_i}$, $m_i=m_{z_i}$, and $I_i=I_{z_i}$. The common
bottom $D$ and lower endpoint $a$ remain those of the original tested
chain; every shortcut and normal-group test is inherited from it.
Write
\[
 T_i=T_i(q_i),\qquad q_i=p_i^{f_i},\qquad n_i=\dim V_i,
 \qquad a_i=|T_i|.
\]
Here $V_i$ is the natural simple module. In characteristic two the
type $B_r$ parameter is $n_i=2r$, using its simple symplectic natural
quotient. For a unitary group the parameter is $q_i=p_i^{f_i}$,
although its natural matrix field is $\mathbb F_{q_i^2}$.
Put $\epsilon_X=2$ for a unitary type and $\epsilon_X=1$ otherwise,
so the natural matrix field has degree $\epsilon_Xf_X$ over its
prime field. We retain this conversion for restriction of scalars;
twisted-tensor steps will have the exact usual-field relation
$f_s=b f_t$.

For each $i<j$ choose the prescribed two-edge shortcut. We write its
ascending steps as $B_{ij},A_{ij}$ and its intermediate simple label as
$Y_{ij}$; a prime denotes an intermediate parameter. The letters $A_{ij}$
and $B_{ij}$ denote steps, not subgroups. On an individual ascending
step, subscripts $s,t$ denote source and target parameters.

All uses of ``arbitrarily long'' below have the same quantifiers: for
every integer $s$ there are such configurations with at least $s$ chain
vertices, and their minimum rank can be required to exceed any prescribed
constant. A finite coloring followed by finite Ramsey preserves this
property. So does retaining every $d$th vertex for a fixed $d$. Constants
in this section are independent of the configuration and of its length.
An explicitly mentioned pool of parameter values may depend on the
configuration, but its cardinality has a uniform bound.

\subsection{Numerical structure of classical maximal steps}

The factor actions needed for the whole-chain reduction were proved in
Proposition~\ref{cls:factor-actions}. We now derive the field, dimension,
and order estimates from the maximal-subgroup descriptions in
Proposition~\ref{cls:literature}. The proofs retain the actual finite
normalizers, since their field actions will enter both the growth
argument and the later constituent calculation.

\begin{proposition}[The maximal-step catalogue]\label{cls:catalogue}
Consider a strict labelled cover and its one-coordinate maximal subgroup
$M$ in an almost simple group with socle $T$, as supplied by
Proposition~\ref{frm:projection}.  Mark its active nonabelian chief factor.
The non-geometric irreducible category below means that its simple
source acts absolutely irreducibly on the ambient natural module
and that the maximal subgroup is outside all geometric collections,
in particular the subfield collection.  All growing-rank assertions
refer to growing \emph{source} rank.
There are absolute constants $C,c>0$ with the following properties.
\begin{enumerate}[label=\textup{(C\arabic*)},start=2]
\item\label{cls:cat-geometric}
For unbounded-rank classical active types, a geometric
defining-characteristic step changes natural dimension by taking a
subspace or quotient, a repeated natural module, a tensor product, or
restriction of scalars.  Subfield steps leave natural dimension
unchanged and multiply the field exponent by a prime.  An
extension-field step, read upwards from the active factor, has
\begin{equation}\label{grw:matrix-field-convention}
 n_t=b n_s,\qquad
 \epsilon_s f_s=b\epsilon_t f_t,
 \qquad f_t=\frac{\epsilon_s}{b\epsilon_t}f_s,
\end{equation}
where $b$ is the prime degree between natural matrix fields.
A twisted-tensor step has
\begin{equation}\label{grw:twisted-field-convention}
 n_t=n_s^b,\qquad f_s=b f_t,
\end{equation}
exactly; here $b$ is prime, after bounded source-rank cases are
removed.  All other changes of field exponent
belong to a fixed finite set of positive rational factors.
\item\label{cls:cat-inactive}
If there is one active simple component and the other geometric
components are soluble, their natural dimensions are absolutely
bounded.  The active dimension in a subspace case is comparable with
the ambient dimension; a totally singular summand and its dual may
represent the same active component twice.  A nontrivial parabolic
radical in these cases has order between $q^{c n_t}$ and $q^{C n_t^2}$.
For a fixed ambient group, once the active slot, other component
dimensions, and the bounded flag/form choices are fixed, there are
boundedly many possible orders of the full inner stabilizer divided by
the order of its active simple factor, with multiplicity included.
\item\label{cls:cat-symplectic}
A symplectic-type step of unbounded active rank is
cross-characteristic.  Its source is symplectic or orthogonal over a
prime field of order $r$, of natural dimension $2d+O(1)$, and its
target natural dimension is $n_t=r^d$.  If $P=M\cap T$, then
\[
 c n_t^2\leq |P|/|T_s|\leq C n_t^2.
\]
Its target field exponent is at most $|T_s|^C$.
\item\label{cls:cat-nonlinear}
After the twisted-tensor case in
\ref{cls:cat-geometric} is separated off, a defining-characteristic almost simple
irreducible step has a single nonzero restricted, nonnatural
Frobenius digit, up to a common Frobenius twist.  This restricted
module may have further tensor factors in a very special
characteristic.  Its dimension is at least $c n_s^2$.  Its source and target field
exponents differ by a rational factor from a fixed finite set.
Natural, dual-natural, and natural-size special-isogeny representations
are included among the geometric cases instead.
\item\label{cls:cat-cross}
For a cross-characteristic almost simple step,
\[
 q_s^{c n_s}\leq n_t\leq |T_s|^C,
 \qquad f_t\leq |T_s|^C.
\]
The upper degree bound also holds in defining characteristic.
For this and the preceding irreducible case, the homomorphism from the
ambient automorphism normalizer of the source simple group to its
automorphism group has kernel of absolute bounded order.  In
particular, a subgroup of that normalizer has order at most
$C$ times the order of its induced automorphism group of the source.
\end{enumerate}
\end{proposition}

\begin{proof}
For the dimension and field statements, take the full stabilizers of
the indicated structures.  On a nondegenerate direct summand the
active group has its natural module; on a totally singular summand
the dual summand has the contragredient module.  A tensor stabilizer
acts as the tensor product of the component natural modules.
Restriction of scalars multiplies dimension by its degree $b$ and
gives $\epsilon_s f_s=b\epsilon_t f_t$ for the natural matrix
fields.  For example, the mixed $\operatorname{GU}_m(q)$ row in
$\operatorname{Sp}_{2m}(q)$ has $b=2$ and equal usual exponents,
since $\epsilon_s=2$ and $\epsilon_t=1$.  It is a restriction of
scalars, with a direct sum after scalar extension, and is not a
twisted-tensor row.  A subfield inclusion changes only the field.
If an ordinary field degree were composite, its proper intermediate
field structure would be preserved by the normalizer, contradicting
maximality.  The classical stabilizer descriptions with prime degree
are given explicitly in
\cite[proofs of Lemmas~5.5 and~5.7]{BurnessLiebeckShalev2017}.
Unitary conjugation and the diagram maps have bounded order, which
accounts for the permitted bounded ratios in the remaining field
operations.  The
twisted-tensor assertion is proved separately in
Lemma~\ref{cls:twisted-tensor} below.

The only soluble classical component groups have bounded natural
dimension: a linear group in growing dimension contains a growing
special linear simple section, and the corresponding assertions for
the form groups follow by their natural hyperbolic subspaces.
Thus the omitted dimensions in a one-active-component subspace or
tensor stabilizer are bounded.  The explicit block description of a
parabolic gives a radical of order $q^e$, with $e\leq Cn_t^2$ and
$e\geq c n_t$ whenever that radical is nontrivial in the retained
large-rank cases.  For example, the linear radical attached to a
$k$-space has exponent $k(n_t-k)$; the symplectic and orthogonal
exponents count the off-diagonal blocks subject to one symmetric or
alternating relation.  The smallest nonzero retained exponent is
linear in $n_t$.  The other factors are the displayed classical
groups and scalar groups.  Consequently their orders are determined
by the dimensions, field, form signs, and the bounded choices of
interchanged flags.  Intersecting the full stabilizer with the ambient
socle adds only its prescribed determinant or spinor condition; it
does not introduce a freely chosen subgroup of a scalar group.  This
proves \ref{cls:cat-inactive}.

For a symplectic-type normalizer put $N=n_t=r^d$, and let $W$ be
the projective image of its normal $r$-group, of order $r^{2d}$.
The full inner-normalizer rows, including their minimal-field
conditions, are given in
\cite[Section~2.2.6, pp.~70--71, Table~2.9]{BrayHoltRoneyDougal2013}.
Their small exceptions disappear on retaining $d\geq4$.
Ordinary linear conjugation fixes scalars pointwise, so it preserves
the nondegenerate scalar commutator pairing
$W\times W\to\mu_r$.  Moreover the projective linear centralizer
of $W$ is $W$: a centralizing operator twists the representing
matrices by a character of $W$, that character is implemented by
an element of $W$ through the pairing, and the remaining centralizer
is scalar by irreducibility.

In the linear and unitary rows, the normal group is an exponent-$r$
extraspecial group for odd $r$, and a central product
$C_4\circ2^{1+2d}$ for $r=2$.  In both cases the ordinary projective
normalizer has quotient $\operatorname{Sp}_{2d}(r)$ on $W$.
Indeed every symplectic change of generators lifts to a center-fixing
automorphism and is implemented by the irreducible representation
with its specified central character.  For $r=2$, multiplying a
generator by the central element of order four restores its required
square; thus a square form attached to a pure extraspecial subgroup
is not an invariant of this scalar-extended normalizer.
In the unitary case an implementing operator preserves the Hermitian
form up to a multiplier, which is removed by a scalar using the
surjective finite-field norm.
The standard tensor generators of $W$ have determinant one: for odd
$r$ this holds for the cyclic shift and diagonal root-of-unity
generators, and for $r=2$ their determinants are
$(-1)^{2^{d-1}}=1$.  The map from this normalizer to the abelian
quotient of the projective linear or unitary group by its simple
socle therefore factors through $\operatorname{Sp}_{2d}(r)$.
That group is perfect in the retained range, so the map is trivial.
Hence the full ordinary projective normalizer is already the inner
intersection $P$.  Quotienting by the entire ambient scalar center
leaves kernel exactly $W$, and
\[
 P/W\cong\operatorname{Sp}_{2d}(r),\qquad
 |P|/|T_s|=(2,r-1)N^2.
\]
This uses the quotient, without requiring the extension to split.

The pure symplectic and orthogonal form rows have $r=2$.
Here the extraspecial square form is invariant.  The cited inner
rows give $P/W=O_{2d}^{\varepsilon}(2)$ or
$\Omega_{2d}^{\varepsilon}(2)$, with active simple group
$T_s=\Omega_{2d}^{\varepsilon}(2)$ and index at most two.
These are the inner rows after the determinant and spinor conditions;
the scalar kernel of the natural cover is removed in $W$.
Consequently $N^2\leq |P|/|T_s|\leq2N^2$ in every retained row.
There is no extra factor from $r-1$ or the ambient field exponent.

The target-field bound uses the minimal-field condition of the
maximal row.  For odd $r$, the degree
$e=\operatorname{ord}_r(p)$ divides $r-1$; the usual target exponent
is $e$ in the linear row and $e/2$ in the unitary row.
The binary and pure form rows have usual exponent one.
Thus $f_t\leq\max(1,r-1)\leq |T_s|$.
The representation characteristic is different from $r$, and the
active source is over the prime field of order $r$.
This proves \ref{cls:cat-symplectic}.

For an absolutely irreducible representation of a central cover
$\widetilde T_s$, the dimension is at most
$|\widetilde T_s|^{1/2}$: the square of each simple-module dimension
is bounded by the dimension of the group algebra.  The field-of-definition bound can be obtained directly, including in
modular characteristic.  For a finite group $K$ in characteristic $p$,
there are at most $|K|$ absolute simple module types, by the semisimple
quotient of its group algebra.  Frobenius permutes these types, so an
absolute simple representation has a Frobenius orbit of length
$a\leq |K|$.  An isomorphism with its $p^a$-Frobenius twist descends the
representation to $\mathbb F_{p^a}$: choose a finite field containing
its matrices and an intertwiner, scale the intertwiner to have norm
one, and apply matrix Hilbert's Theorem~90.  The scalar norm adjustment
is possible because finite-field norms are surjective.  Apply this
with $K=\widetilde T_s$.  The Schur multipliers of the
large classical groups have order at most a polynomial in $|T_s|$;
the finitely many exceptional multipliers do not change the bound.
Forms cost an extension of degree at most two.  Maximality outside
the subfield collection makes these bounds applicable to the target
field itself.  For the cross-characteristic lower bound, the
classical tables in \cite[Section~3, Tables~3--6]{Hasa2014} give
degree at least $q_s^r/3$ once the source rank $r\geq6$, including
the smaller Weil degrees.  Their excluded families have bounded
rank.  As $n_s\leq2r+1$, this is at least $q_s^{n_s/4}$ after
increasing the rank threshold.  This proves the numerical assertions
of \ref{cls:cat-cross}.

For the defining-characteristic quadratic bound, use
\cite[Theorem~5.1 and Table~2]{Luebeck2001} in rank $r>11$.
Modules outside that table have dimension greater than $r^3/8$
in type $A$, and greater than $r^3$ in types $B,C,D$.
Every listed nontrivial, nonnatural, nondual degree is at least $r^2/2$.
Consequently the retained restricted module has dimension at
least $r^2/2\geq n_s^2/10$.  Type $B$ in characteristic two is
read through its type-$C$ natural simple quotient, as in our
parameter convention; the natural-size special-isogeny cases are
already assigned to geometry.

A single occupied ordinary Frobenius digit has full source period:
any shift strictly between zero and $f_s$ moves it to a different
position in the unique finite restricted label.  Returning through
$f_s$ applies at most the bounded diagram operation.  An internal
tensor decomposition of that digit changes neither its occupied
position nor this uniqueness argument.  Thus the minimal scalar
field degree is $f_s$ times one of the bounded diagram periods.
The matrix degree is separately $\epsilon_t f_t$.  The
form-compatible minimal-field argument in
Lemma~\ref{cls:normalizer} bounds the ratio of this matrix degree
to the minimal degree outside the geometric collections.
This gives the fixed finite rational set of usual exponent ratios
and proves \ref{cls:cat-nonlinear}.
Lemma~\ref{cls:normalizer} below supplies
the final normalizer assertion.
\end{proof}

\begin{lemma}[Irreducible normalizers]\label{cls:normalizer}
In a non-geometric absolutely irreducible step of
Proposition~\ref{cls:catalogue}, let $Y$ denote the simple source
inside the ambient projective classical group $T$.  In particular,
the maximal subgroup is not in a subfield collection.  Then
\[
 \bigl|\ker\bigl(N_{\operatorname{Aut}(T)}(Y)
       \longrightarrow\operatorname{Aut}(Y)\bigr)\bigr|\leq C
\]
for an absolute $C$.  Moreover
$|\operatorname{Out}(Y)|\leq C n_Y f_Y$ for a large-rank classical
source with parameters $(n_Y,p^{f_Y})$.
\end{lemma}
\begin{proof}
Lift the irreducible projective representation to the perfect central
cover of $Y$.  A projective linear element centralizing its image
commutes with that representation up to a linear character of the
perfect cover, hence commutes exactly.  Schur's Lemma makes it scalar,
so it disappears projectively.

If a semilinear element in this kernel has scalar-field action $\sigma$,
the representations $\rho^\sigma$ and $\rho$ are isomorphic.  If
$\sigma$ has order $a$ on the scalar field $E$, its intertwiner $A$
satisfies
\[
 A\sigma(A)\cdots\sigma^{a-1}(A)=cI,
 \qquad c\in E^{\sigma,*}.
\]
Surjectivity of the scalar norm adjusts $A$ so that the product is
one; matrix Hilbert's Theorem~90 then gives a model over $E^\sigma$.
An invariant bilinear form is a line defined by linear equations
over the minimal field, and invariant quadratic-form coefficients
satisfy linear equations there as well.  In characteristic two,
two invariant quadratic forms with the same polarization differ
by the square of an invariant linear form.  Absolute irreducibility
makes that difference zero, so the prescribed quadratic form itself
descends up to scalar.  Matching finite-field form classes costs at
most a quadratic extension.  For a unitary group the involution on
the minimal field is either nontrivial, giving Hermitian descent
and surjective norms, or trivial, giving a preserved bilinear
structure and hence a geometric form case.  A larger field period would
therefore provide a proper subfield structure normalized by the
whole irreducible normalizer, contrary to this step being outside the
subfield collection.

To justify the assertion about the whole
normalizer, put an absolutely irreducible model over its minimal
field $E_0$.  A normalizing matrix intertwines this model with the
model obtained by a source automorphism and a field automorphism;
both have entries in $E_0$.  The intertwining equations therefore
have coefficients in $E_0$.  Their solution space over the algebraic
closure is one-dimensional by absolute irreducibility.  It consequently
has a nonzero $E_0$-rational solution, necessarily invertible.  Up to a
scalar, the normalizer thus preserves the projective $E_0$-structure.
The same argument applies after dualizing in a graph component, and
the uniqueness of the invariant form gives the already allowed
quadratic extension.  A graph operation contributes at most another
factor of two in unbounded rank.  This bounds the kernel independently
of the ambient field degree.

Finally the diagonal, field, and graph
description of automorphisms gives respectively factors at most
$n_Y$, $f_Y$, and an absolute constant in the outer group.
\end{proof}

Twisted tensor subgroups and their normalizers were studied in detail by
Schaffer \cite{Schaffer1999}; the algebraic and finite tensor structures
used here are those of \cite[Section~4(5)]{LiebeckSeitzClassical1998}.
The next lemma proves the precise maximality consequence and field relation
needed for these tested chains.

\begin{lemma}[Twisted tensor maximality]\label{cls:twisted-tensor}
An almost simple irreducible maximal subgroup with a nontrivial
Steinberg tensor decomposition is, in growing source rank, a geometric
twisted-tensor step with the exact dimension and field relations in
\ref{cls:cat-geometric}.  Its tensor-orbit degree is prime.
\end{lemma}
\begin{proof}
Let $S$ be the simple source, $M$ the almost simple maximal subgroup,
and $T$ the ambient simple classical group.  Work over an algebraic
closure on the cover supplied by Lemma~\ref{lift:finite-covers}, and
write the irreducible module as
\[
 V=\bigotimes_{i\in I}V_i,
 \qquad
 \mathcal A_i=\operatorname{End}(V_i)\otimes1
       \ \leq\operatorname{End}(V),
\]
where $V_i$ are the nonzero ordinary restricted Frobenius digits,
with their indicated twists.  The dimensions of these nontrivial
digits tend to infinity with the source rank.

We first verify preservation of the actual tensor structure.  A
source field automorphism translates ordinary digit positions,
applying the diagram operation when a position passes the source
period; inner, diagonal, and diagram automorphisms act on the
restricted modules at those positions.  Ambient semilinear
operations give the same shifts, and a graph operation dualizes
the modules.  Uniqueness of the finite restricted highest-weight
parametrization matches the resulting whole digit modules.  Tensor
their individual intertwiners, with the matching permutation, to
obtain an intertwiner of $V$.  The discrepancy between a projective
and linear intertwining identity is a linear character of the
perfect cover and is trivial.  Absolute irreducibility makes the
space of whole-module intertwiners one-dimensional.  Every actual
normalizing operator therefore differs from this tensor intertwiner
by a scalar and permutes the embedded algebras $\mathcal A_i$.
For a graph operation the corresponding map on endomorphisms is an
anti-isomorphism, which has the same conclusion.  This argument
preserves a whole ordinary digit even if that restricted module
has an internal tensor decomposition in a very special
characteristic.

Replace each digit by its full connected special linear group, or
by its full connected bilinear-form group when required.  Bilinear
self-duality makes each ordinary digit self-dual separately.
Unitary descent acts on the special linear factors, permitting
duals.  In characteristic two every nontrivial irreducible self-dual
digit has an alternating form: a nonalternating invariant form
would give an invariant linear functional by taking square roots
of its diagonal.  The tensor product of at least two nondegenerate
alternating spaces has a quadratic form with the product
polarization and value zero on pure tensors
\cite[Proposition~2.2(iv)]{LiebeckSeitzClassical1998}.
It agrees with the prescribed finite quadratic form, after a scalar
normalization of the polarization: their difference is the square
of an invariant linear form, which vanishes by nontrivial absolute
irreducibility.  These constructions give a connected semisimple
tensor subgroup $K$ in the ambient form group.

The groups are intrinsic to the algebras $\mathcal A_i$ and their
restricted adjoint involutions.  Hence the actual ambient Frobenius
$F$ and the full finite normalizer preserve $K$.  In growing
classical rank, finite ambient automorphisms have lifts commuting
with $F$; diagonal lifts satisfy this projectively because
$F(x)x^{-1}$ is central.  Graph and field lifts also commute, with
only the bounded-rank triality and exceptional graph cases removed.
This is the setup of
\cite[Theorem~2 and Section~4]{LiebeckSeitzClassical1998}.
Henceforth view $K$, and later $K'$, in the ambient adjoint group.  Thus $M$ normalizes the terminal derived subgroup
$H=(K^F)^{(\infty)}$.  Since $H$ is perfect it lies in $T$, and
$S\leq H$ because $S$ is a perfect subgroup of $K^F$.

For each orbit of $F$ on the factors, fixed points are obtained by
choosing one coordinate fixed by the corresponding power of $F$;
see \cite[Proposition~2.8(ii)]{LiebeckSeitzClassical1998}.
Consequently $H$ is a central product of classical components, one
per orbit.  It is proper in $T$ also as a finite group.  To check
this uniformly, put $q=p^{f_t}$, let the individual digit dimensions
be $d_1,\ldots,d_k$, put $d_{\min}=\min_i d_i$, and write
$n_t=\prod_i d_i$.  The earlier bound for nontrivial digits gives
$d_{\min}\to\infty$ with the source rank.
Elementary classical order bounds give
\[
 |H|\leq q^{2\sum_i d_i^2},\qquad
 |T|\geq q^{n_t^2/8}.
\]
For $k\geq2$,
$\sum_i d_i^2/n_t^2\leq k/d_{\min}^{2(k-1)}\leq2/d_{\min}^2$.
Thus $H<T$ for all sufficiently large source ranks.  If $MH$
were the entire almost simple ambient group, its normal subgroup
$H\leq T$ would equal $T$.  Hence $MH$ is proper, and maximality
gives $H\leq M$.  But the perfect group $H/S$ is a subgroup of the
soluble group $M/S$, so $H=S$.  There is therefore one Frobenius
orbit, and its full natural classical component is the source.
If its length is $b$, the digit dimension is $n_s$ and
$n_t=n_s^b$ exactly.

The field relation is read from this actual natural-factor cycle.
Put $t=f_t$.  For a linear target, $F^b$ on one factor is the
standard field map with exponent $bt$, up to an inner conjugacy.
For a unitary target, write $F=F_{p^t}\tau$ with $\tau$ the
inverse-transpose operation.  On one factor the return map is
$F_{p^{bt}}\tau^b$, again up to an inner conjugacy.  Its source is
unitary when $b$ is odd and linear when $b$ is even, in both cases
with usual exponent $bt$.  An orthogonal graph choice changes the
form sign, not this exponent.  The characteristic-two quadratic
construction likewise changes no field exponent.  Thus
\begin{equation}\label{cls:tensor-field}
 f_s=b f_t
\end{equation}
exactly, in the convention of this paper.  The natural factors have
positions $0,t,\ldots,(b-1)t$ modulo $bt$, which are distinct.
For the odd unitary convention using positions $2jt$, multiplication
by two permutes the same residues modulo $b$; the prime two-cycle
in a unitary target has a linear source.  These are precisely the
usual-field alternatives in
\cite[Section~4(5)]{LiebeckSeitzClassical1998}.

Finally suppose $b=uv$, with $u,v>1$.  Number the factors by
$\mathbb Z/b\mathbb Z$, and group the cosets of its unique subgroup
of order $u$.  The normalizer's permutation normalizes the regular
cyclic descent action and preserves this block system.  The
algebra for a block $B$ is the algebra generated by its
$\mathcal A_i$, namely $\operatorname{End}(\bigotimes_{i\in B}V_i)
\otimes1$.  Replace each block by its full special linear or
bilinear-form group.  For a characteristic-two quadratic target
use the full symplectic group on each grouped alternating space.
The resulting quadratic tensor form is unchanged: it has the
same polarization and vanishes on every fully pure original
tensor, and such tensors span $V$.  Hence this grouped subgroup
$K'$ still preserves the prescribed ambient form and is invariant
under $F$ and $M$.

Its terminal derived fixed-point subgroup $H'$ has one classical
component on dimension $n_s^u$ over the field with usual exponent
$v f_t$, with the unitary parity or orthogonal sign prescribed by
the return map.  For $n_s\geq16$, the elementary bounds
$Q^{d^2/8}\leq|\operatorname{Cl}_d(Q)|\leq Q^{2d^2}$ give
\[
 |S|\leq q^{2uv n_s^2}
   <q^{v n_s^{2u}/8}\leq|H'|,
 \qquad
 |H'|\leq q^{2v n_s^{2u}}
   <q^{n_s^{2uv}/8}\leq|T|.
\]
The strict inequalities follow from
$n_s^{2u-2}>16u$ and $n_s^{2u(v-1)}>16v$.
As before $M$ normalizes $H'$ and $S\leq H'$.  Its properness in
$T$ forces $H'\leq M$, and perfectness then forces $H'=S$,
contradicting the first inequality.  Therefore $b$ is prime.
\end{proof}

\subsection{Estimates along a two-edge shortcut}

We record how the catalogue applies to the reduced configurations.
In the estimates below, a \emph{nonlinear} cover means a
non-geometric defining-characteristic irreducible cover as in
\ref{cls:cat-nonlinear}; a \emph{crossing} cover means an almost simple
cross-characteristic cover or a symplectic-type cover.
The field estimates constrain prime degrees, the dimension estimates
can be iterated along the chain, and the inner order quotients will
multiply through intermediate vertices.

\begin{enumerate}[label=\textup{(G\arabic*)},ref=G\arabic*]
\item\label{grw:input-fields}
For a geometric or defining-characteristic almost simple cover,
\ref{cls:cat-geometric} and \ref{cls:cat-nonlinear} express the field
ratio as a fixed rational factor times the contribution of a subfield,
extension-field, or twisted-tensor degree. An exceptional degree is
prime or belongs to a fixed bounded set of positive integers.
All bounded rational corrections, including two, are fixed before
any prime coloring is made. The restriction-of-scalars and
twisted-tensor formulas are \eqref{grw:matrix-field-convention}
and \eqref{grw:twisted-field-convention}, respectively.
The mixed unitary restriction of degree two is a bounded geometric
case, with $\epsilon_s/\epsilon_t\in\{1/2,1,2\}$.
A natural space and its dual give a separate dimension operation,
not an additional restriction degree. Other form and special-isogeny
rows retain their own fixed rational field factors; no telescoping
of $\epsilon$ factors through those rows is used.
Lemma~\ref{cls:twisted-tensor} supplies distinct tensor positions
modulo $f_s$, including the unitary alternatives. Thus there is no
matrix-field factor of two in \eqref{grw:twisted-field-convention},
and the dimension formulas for a two-step subfield/extension or
subfield/twisted-tensor path are exact.

\item\label{grw:input-geometric}
The geometric dimension operations in \ref{cls:cat-geometric},
together with Proposition~\ref{cls:reduction}(v), give boundedly
many omitted dimensions once the active slot, permutation data,
and inactive simple labels are fixed. In paths with no replication
and no inactive nonabelian factor, these dimensions are bounded
integers; ordinary tensor factors other than the active one then
have bounded dimensions. An unbounded replication degree belongs
to the bounded-cardinality pool in Proposition~\ref{cls:reduction}.

\item\label{grw:input-degrees}
Clause~\ref{cls:cat-nonlinear} gives $n'\ge c n^2$ on a nonlinear
cover, with nontrivial Steinberg tensor decompositions already
assigned to geometry. Clauses~\ref{cls:cat-symplectic} and
\ref{cls:cat-cross} give the common crossing estimate
\begin{equation}\label{grw:cross-bounds}
 n'\ge \exp(c n\log q),\qquad n',f'\le |T|^C.
\end{equation}
For an almost simple cover the upper degree bound also holds in
defining characteristic.
Every ascending same-characteristic geometric step decreases natural
dimension by at most a fixed factor, by
\ref{grw:input-fields}--\ref{grw:input-geometric}. It also decreases
$n\log q$ by at most a fixed factor. For restriction of scalars,
\eqref{grw:matrix-field-convention} gives
\[
 n_t\log q_t
 =\frac{\epsilon_s}{\epsilon_t}n_s\log q_s
 \ge\tfrac12 n_s\log q_s,
 \qquad n_t=b n_s\ge n_s.
\]
For a twisted-tensor step,
\eqref{grw:twisted-field-convention} gives
$n_t\log q_t=(n_s^b/b)\log q_s\ge n_s\log q_s$ when $n_s,b\ge2$.
The other geometric rows have fixed positive lower factors. A
same-characteristic nonlinear step has both lower-bound properties
as well, by its quadratic degree bound and fixed positive field ratio.

\item\label{grw:input-multipliers}
In a one-factor almost simple cover the inner order quotient is at
most $|\operatorname{Out}(T)|$. Clause~\ref{cls:cat-symplectic}
bounds the symplectic-type quotient between $cN^2$ and $CN^2$,
where $N$ is the target natural dimension; the target dimension
and field exponent are at most $|T|^C$ for source $T$.
For a parabolic or unipotent subspace cover with one active factor,
no inactive nonabelian factor, and bounded field ratio,
\ref{cls:cat-inactive} makes the target dimension $n'$ comparable
to the source dimension $n$ and places the inner quotient between
$q^{c n'}$ and $q^{C n^2}$. For fixed target type and parameters
that quotient has boundedly many possibilities. This is a bound
on cardinality; the quotient itself is usually very large.
\end{enumerate}

A strict split followed by an equal-type merge is geometric.
Every shortcut containing an almost simple irreducible or a
symplectic-type cover has one active simple factor at both covers
and no other nonabelian inner
composition factor, by Proposition~\ref{cls:reduction}. For those
shortcuts, the next lemma transfers the normalizer bound to the
actual lower coordinate stabilizer.

\begin{lemma}[Transport of a coordinate stabilizer]\label{grw:input-normalizers}
Suppose that the last cover is an almost simple irreducible cover with
source $Y$, and that all three shortcut multiplicities agree, as in
every application below.  Denote the lower node's full upper-coordinate
stabilizer image by $M_{ij}$, and let $P_B$ be the inner projection in
the first cover.  Then
\begin{equation}\label{grw:normalizer-input}
 |M_{ij}|\le C\,|P_B|\,n'^{C}(1+f')^{C}.
\end{equation}
\end{lemma}
\begin{proof}
 Write the actual shortcut as
$U=z_i\prec V\prec W=z_j$, and work modulo $F_W$.
In one $W$-coordinate the last cover has
$Y\trianglelefteq P_A\le\operatorname{Aut}(Y)$.  Its killed kernel
from the middle is soluble by Proposition~\ref{cls:reduction} and
normal in $P_A$, hence is trivial.  The active subgroup in that
coordinate is therefore the actual simple group $Y$.
Proposition~\ref{frm:active-transport} gives a $V$-equivariant bijection
$\beta:I_V\longrightarrow I_W$ because the multiplicities agree.
For $\alpha\in I_V$ and $t=\beta(\alpha)$, the actual stabilizers in
$U$ coincide: $U_\alpha=U_t$.  Let $\rho_V,\rho_W$ be their middle
and upper coordinate actions, and let
$\phi:N_{\operatorname{Aut}(T_j)}(Y)\longrightarrow\operatorname{Aut}(Y)$
be conjugation.  The quotient-induced actions satisfy
\[
 \phi\circ\rho_W|_{U_t}=\rho_V|_{U_\alpha},
 \qquad \phi(M_{ij})=M_B,
\]
where $M_B$ is the actual first-cover coordinate stabilizer image.

The inner-intersection assertions use the original lower endpoint $a$
from Definition~\ref{cls:tested-chain}.  All retained comparisons
start strictly above this $a$, so none uses an excluded middle
$v_{0j}$.  Condition
\ref{cls:test-alternating} makes $a$ an intersection of coatoms in
$[a,V]$ and $[a,W]$ for every retained middle and upper endpoint.
Proposition~\ref{frm:projection} in these original-$a$ projections
identifies the inner part of a saturated image with its transported
inner projection.  For $U\prec V$ the coatom case of that Proposition
also applies directly.  In particular,
\[
 M_B\cap\operatorname{Inn}(Y)=P_B.
\]
There is no assumption that a meet of coatoms in the original global
interval is a meet of coatoms in every smaller interval.
The bounded kernel from Lemma~\ref{cls:normalizer} now gives
\[
 |M_{ij}|\le C|M_B|
 \le C|P_B|\,|\operatorname{Out}(Y)|
 \le C|P_B|n'f',
\]
which implies \eqref{grw:normalizer-input}.  This includes diagonal,
field, and graph operations.  All higher-core losses are already
kernels of the displayed actual coordinate actions; no uncontrolled
factor from the field of the final target is inserted.
\end{proof}

\subsection{Finite colors for field degrees}

For a positive rational $x$, let $v_\ell(x)$ be its valuation at a prime
$\ell$.  If $S$ is a finite set of primes, put
\[
 \Omega_S^+(x)=\sum_{\ell\notin S}\max\{v_\ell(x),0\},
 \qquad
 \Omega_S^-(x)=\sum_{\ell\notin S}\max\{-v_\ell(x),0\}.
\]
Both sums have finite support.

\begin{lemma}[Finite prime complexity]\label{grw:prime-colors}
Let $x_1,\ldots,x_s$ be positive real numbers whose pair ratios are
positive rationals.  Suppose that, for every $i<j$,
\[
 x_j/x_i=c_{ij}u_{ij}^{\varepsilon},\qquad
 c_{ij}\in\mathcal C,\qquad \varepsilon\in\{1,-1\},
\]
where $\mathcal C$ is a fixed finite set of positive rationals,
$u_{ij}$ is a positive integer with at most $h$ prime factors counted
with multiplicity, and $\varepsilon$ is the same for all pairs.
There is a number of colors depending only on $\mathcal C,h$ such that
on every homogeneous triangle all three $x$'s are equal.
Consequently arbitrarily long such arrays have arbitrarily long
constant subarrays.
\end{lemma}

\begin{proof}
Let $S$ contain the primes occurring in the numerators or denominators
of members of $\mathcal C$.  For $\ell\in S$, the integer
$v_\ell(x_j/x_i)$ lies in a fixed bounded interval: the contribution of
$c_{ij}$ is bounded, and that of $u_{ij}$ has absolute value at most $h$.
Outside $S$ all valuations have sign $\varepsilon$, and the sum of their
absolute values is at most $h$.  Color $(i,j)$ by the valuation vector
on $S$ and this last sum.  This is a finite coloring with the asserted
uniformity.

For $i<j<k$ we have
\[
 v_\ell(x_k/x_i)=v_\ell(x_j/x_i)+v_\ell(x_k/x_j).
\]
On a homogeneous triangle each coordinate of the common valuation
vector therefore equals twice itself, so it is zero.  Outside $S$
there is no cancellation because all valuations have the same sign.
Their total absolute value likewise equals twice itself, hence is zero.
Every valuation of every pair ratio is zero, so all ratios equal one.
Finite Ramsey gives the last assertion.
\end{proof}

\begin{lemma}[Opposite field moves]\label{grw:opposite-fields}
Arbitrarily long homogeneous shortcut chains cannot have, on every
comparison, one proper subfield step and one proper extension step,
including the twisted-tensor extension alternative.
\end{lemma}

\begin{proof}
Color first by the order and kind of the two steps.  By
\ref{grw:input-fields}, their combined dimension formula is either
\[
 n_j=b_{ij}n_i\quad\hbox{for every }i<j,
 \qquad\hbox{or}\qquad
 n_j=n_i^{b_{ij}}\quad\hbox{for every }i<j,
\]
where $b_{ij}>1$ is prime or belongs to a fixed bounded set of integers.
In the first case put $x_i=n_i$, and in the second put $x_i=\log n_i$.
In either case $x_j/x_i=b_{ij}$ has uniformly bounded prime complexity.
Lemma~\ref{grw:prime-colors}, with $\mathcal C=\{1\}$, gives a
homogeneous triangle with $b_{ij}=1$, a contradiction.  Notice that
field ratios alone would not suffice: ratios $r_j/r_i$ of distinct
primes are compatible on arbitrarily many vertices.  It is the exact
dimension formula that excludes this possibility.
\end{proof}

\begin{proposition}[Purely geometric paths]\label{grw:no-geometric}
An arbitrarily long reduced chain cannot have exclusively geometric
two-edge shortcuts.  This includes strict split/equal-type merge
shortcuts.
\end{proposition}

\begin{proof}
Color by the two cover categories, the endpoint family, and whether a
split and merge occur.  There are finitely many colors.  The endpoint
characteristic is common because every step is in defining
characteristic.  If both an increasing and a decreasing exceptional
field move occur, Lemma~\ref{grw:opposite-fields} applies.  Otherwise
all exceptional primes in a pair ratio have the same sign.  The product
of the two bounded rational factors belongs to a fixed finite set, and
at most two exceptional degrees occur.  A bounded exceptional degree
has uniformly bounded prime complexity as well.  Apply
Lemma~\ref{grw:prime-colors} to $x_i=f_i$.  We may therefore assume
that all endpoint fields and families agree.

The exceptional degrees at the individual covers are now bounded
integers.  Indeed, with no opposite field moves their product, or its
reciprocal, is a member of the fixed finite set of bounded twisting
ratios, because $f_j/f_i=1$.  Each positive degree divides the numerator
or denominator of one of these finitely many rationals.

Fix a lower endpoint $i$.  We show that its upper natural dimension has
only boundedly many possible values.  For block and parabolic shapes,
the ambient dimension is the active dimension, or twice that dimension
when its dual is present, plus the dimensions of the omitted blocks.
The dimensions of nonabelian omitted blocks are determined, with bounded
ambiguity, by the bounded-cardinality pool of inactive types.  Soluble
omitted blocks have bounded dimensions.  A tensor shape multiplies the
active dimension by a product of the omitted dimensions.  A replicated
base multiplies it by a replication degree from its bounded-cardinality
pool.  A form change has the same natural dimension, with the bounded
natural-module exceptions included in the family color.  Field
restriction multiplies dimension by one of the bounded degrees just
obtained, and a twisted-tensor step raises it to one of those degrees.
Thus one geometric cover gives a bounded list of possible outputs from
a fixed input.  Two covers give a list whose size is at most the product
of the two bounds.  For a split followed by a merge the second step has
the same simple label, so gives no new natural dimension at all.

All labels now have the same family and field.  Their strict growth in
order makes their natural dimensions distinct.  There cannot be more
upper endpoints than the bounded number of possible outputs from the
fixed lower endpoint.  This contradicts arbitrary length.
\end{proof}

\subsection{Multiplicative inner quotients}

From now on each comparison contains a nonlinear or crossing step.
There is no splitting.  The multiplicity is the common integer $m$
from Proposition~\ref{cls:reduction}, also at the two covers of every
chosen shortcut.  Let $P_{ij}\le T_j$ be the saturated inner projection
of the lower node at an upper coordinate, and let
$M_{ij}\le\operatorname{Aut}(T_j)$ be its full induced coordinate
stabilizer.  Applying Proposition~\ref{frm:projection} with the
original $a$ and condition~\ref{cls:test-alternating}, as explained
in Lemma~\ref{grw:input-normalizers}, gives
$P_{ij}=M_{ij}\cap T_j$.

\begin{lemma}[Index quotients]\label{grw:indices}
The positive integers
\[
 e_{ij}=|P_{ij}|/|T_i|
\]
satisfy $e_{ik}=e_{ij}e_{jk}$ for every $i<j<k$.  They multiply in the
same way through each selected shortcut.
\end{lemma}

\begin{proof}
The active factor $T_i$ is a composition factor of $P_{ij}$, so its
order divides $|P_{ij}|$.  Saturation and supplementation give
\[
 [z_j:z_i]=[T_j^m:P_{ij}^m]
 =\left(\frac{|T_j|}{|P_{ij}|}\right)^m.
\]
Multiplying group indices through $z_j$ and taking the positive $m$th
root gives
\[
 \frac{|T_k|}{|P_{ik}|}
 =\frac{|T_k|}{|P_{jk}|}\frac{|T_j|}{|P_{ij}|}.
\]
Rearranging and dividing by $|T_i|$ proves the formula.  The same
calculation applies at a shortcut because its three multiplicities
agree.  No identification of its intermediate label with a chain label
is required.
\end{proof}

We use the soluble-normal and commuting-normal conclusions of
Lemma~\ref{frm:labels} and Corollary~\ref{frm:bottleneck}.
Fix retained labelled endpoints $z_i<z_j$ and put
$F_j=\operatorname{core}_{z_j}(D)$.  This original higher core is
normal in $z_j$ and contained in $D$, hence normal in every
intervening group.  Write
\[
 S_j=N_j/F_j=\prod_{t\in I_j}T_{j,t},\qquad
 (z_i\cap N_j)/F_j=P_{ij}^{I_j}.
\]
For each $M_{ij}$-invariant subgroup $R\le T_j$, write
$\widehat R\le S_j$ for its transported coordinate product.
It is normalized by $z_i/F_j$.

\begin{lemma}[Coordinate normal tests]
\label{grw:coordinate-normal-tests}
Let
\[
 z_i/F_j\le\overline U<z_j/F_j
\]
be an actual subgroup, and let $U$ be its full inverse image in
$z_j$.  A transported product normal in $\overline U$ is called
\emph{killed} if it is contained in $D/F_j$; it then lies in
$\operatorname{core}_{\overline U}(D/F_j)$.
\begin{enumerate}[label=\textup{(N\arabic*)},ref=N\arabic*]
\item\label{grw:normal-soluble}
If $\widehat R\trianglelefteq\overline U$ is soluble, then
\begin{equation}\label{grw:normal-containment}
 \widehat R\le D/F_j,\qquad R\le P_{ij}.
\end{equation}
\item\label{grw:normal-commuting}
Let $R,Q\le T_j$ be $M_{ij}$-invariant, with
$\widehat R,\widehat Q$ normal in the same $\overline U$.
Suppose that
\begin{equation}\label{grw:commuting-kernel}
 [\widehat R,\widehat Q]\le\overline K
       \trianglelefteq\overline U,\qquad
 \overline K\le D/F_j.
\end{equation}
Then at least one of $\widehat R,\widehat Q$ is killed, and
its coordinate subgroup is contained in $P_{ij}$.
\end{enumerate}
\end{lemma}

\begin{proof}
For $U>z_i$, the original test \ref{cls:test-normal} provides
a fenced interval or a labelled bottleneck at $U$.  If $U=z_i$,
the endpoint itself is labelled.  Since $F_j\trianglelefteq z_j$
and $F_j\le D$, apply Lemma~\ref{frm:labels} or
Corollary~\ref{frm:bottleneck} in the quotient by $F_j$, with
bottom $D/F_j$.  The soluble-normal conclusion gives
\eqref{grw:normal-containment}.  For
\textup{(\ref{grw:normal-commuting})}, the full inverse image of
$\overline K$ is normal in this same $U$ and contained in $D$,
so the commuting-normal conclusion kills one product.
Since $D\le z_i$, a killed product has its coordinate subgroup
in $P_{ij}$.
\end{proof}

For \textup{(\ref{grw:normal-soluble})} we use the enlargement
\begin{equation}\label{grw:normal-enlargement}
 \overline U=(z_i/F_j)\widehat R;
\end{equation}
for \textup{(\ref{grw:normal-commuting})} we use
$\overline U=\langle z_i/F_j,\widehat R,\widehat Q\rangle$.
Properness and normality in this same enlargement are checked
in each application.  If one tested product contains a transported
active subgroup known to map onto the original lower labelled
socle, it is not killed, so the other product is killed.
Containment in $P_{ij}$ alone does not imply that a product
is killed and cannot replace the normality and bottom-containment
hypotheses on $\overline K$.  Solubility here is a property of
the quotient product; its full inverse image need not be soluble.
No normality of the lower core $F_i$ in the enlargement, or full
supplementation dichotomy at an unlabelled bottleneck, is used.

\Needspace{5\baselineskip}
\subsection{A small projective group and a parabolic}

\begin{lemma}[Small coordinate stabilizers]\label{grw:small-parabolic}
There is an absolute constant $c_0>0$ with the following property.
Let $T$ be a classical simple group of sufficiently large natural
dimension $n$, over parameter $q=p^f$, and let
$M\le\operatorname{Aut}(T)$ have order $m$.  If
\[
 n>2(m^2+m),
\]
then $M$ normalizes a proper parabolic subgroup of $T$ whose unipotent
radical has order at least $q^{c_0 n}$.  Consequently the coordinate
normal test excludes $M=M_{ij}$ whenever, in addition,
$q_j^{c_0 n_j}>|M_{ij}|$.
\end{lemma}

\begin{proof}
Use the natural-space field $E=\mathbb F_{p^e}$; thus $e=f$, or $e=2f$
in the unitary case.  Apart from the bounded-rank graph exceptions
already excluded, automorphisms in the form cases are represented by
projective semilinear similarities.  In the linear case there may also
be correlations, which interchange the natural space and its dual.
The ordinary orthogonal graph operation is a similarity in this
description; characteristic-two type $B$ uses its symplectic natural
quotient.

First suppose that $M$ acts on the natural-space side.  Choose one
semilinear representative $s_g$ for each projective element, with
$s_1=1$.  The products satisfy
$s_gs_h=c(g,h)s_{gh}$ for nonzero scalars $c(g,h)\in E$.
For an unknown vector $v\in E^n$ impose
\[
 B(s_gv,s_hv)=0\quad(g,h\in M)
\]
in a bilinear or Hermitian form case, and impose also
$Q(s_gv)=0$ in a quadratic form case.  In a plain linear space no
equations are necessary.  Field automorphisms and semilinear maps are
$\mathbb F_p$-linear.  These are therefore at most $e(m^2+m)$ scalar
polynomial equations of degree at most two in $en$ prime-field
variables.  Their sum of degrees is less than $en$.

For completeness, the elementary polynomial count giving a nonzero
solution is as follows.  If $F_1,\ldots,F_t$ have total degree sum less
than the number $N$ of prime-field variables, the number of their common
zeros, reduced modulo $p$, is
\[
 \sum_{x\in\mathbb F_p^N}
       \prod_{a=1}^t(1-F_a(x)^{p-1}).
\]
Every expanded monomial has degree less than $N(p-1)$, so some variable
has exponent less than $p-1$.  The sum over that variable is zero in
$\mathbb F_p$, including exponent zero.  The displayed sum is zero.
All our equations vanish at zero, so their common zero set has another
element.  This is the needed Chevalley--Warning argument
\cite{Chevalley1935,Warning1935}.

Choose such a nonzero $v$ and put
\[
 U=\operatorname{span}_E\{s_gv:g\in M\}.
\]
Then $1\le\dim U\le m$.  Scalar discrepancies between representatives
do not change this span, so $U$ is $M$-stable.  The imposed equations
make $U$ totally singular in the form cases.  In particular, in
characteristic two we used the quadratic equations as well as the
polar equations.  Thus the stabilizer of $U$ is a proper parabolic
normalized by $M$.

Now let $T$ be linear and suppose $M$ contains correlations.  Put
$H=M\cap\operatorname{P\Gamma L}(V)$ and choose $d\in M\setminus H$.
On subspaces the correlation has the form
\[
 d(S)=\operatorname{ann}(F(S))
\]
for an invertible semilinear map $F:V\longrightarrow V^*$.
Choose representatives $s_h$, $h\in H$, and impose
\[
 F(s_hv)(s_kv)=0\qquad(h,k\in H).
\]
Again there are at most $e|H|^2$ quadratic scalar equations.  The same
count gives $v\ne0$.  Let
$U=\operatorname{span}_E\{s_hv:h\in H\}$ and $W=d(U)$.
The equations say $U\le W$.  Since $H$ is normal in $M$, both subspaces
are $H$-stable.  Moreover $d(W)=d^2(U)=U$, because $d^2\in H$.
Thus the flag $U<W$ is preserved by all of $M$, as an unordered pair
under correlations.  Its dimensions are $r,n-r$, where
$1\le r\le |H|<n/2$, so it is a proper flag.  Its parabolic stabilizer
is normalized by $M$.  No involution assumption on $d$ has been used.

In these constructions the small subspace has dimension
$1\le r\le m<n/4$ once $n$ is sufficiently large.  The upper triangular
block matrices in the radical supply at least $n/4$ independent
prime-parameter root coordinates, each of order at least $q$.
For example the linear case has $r(n-r)$ coordinates, and the flag
case has at least $r(n-2r)$; in the polar cases the block from the
nondegenerate middle space to the singular $r$-space supplies at least
$r(n-2r-2)$ coordinates.  These lower bounds are at least $n/4$ after
enlarging the fixed dimension threshold.  We may take a smaller
absolute $c_0$ to cover all the families.  Projectivization does not
remove unipotent coordinates: the scalar center has order prime to $p$.

Let $J$ be this proper parabolic and $R$ its radical.  The group $R$
is soluble and normalized by $M$.  In the coordinate application,
take the actual enlargement \eqref{grw:normal-enlargement}.  Its inner
part is
\[
 \overline U\cap S_j=(P_{ij}R)^{I_j},\qquad
 P_{ij}R\le N_{T_j}(J)<T_j.
\]
The last inequality holds because $J$ is a nontrivial proper subgroup
of the simple group $T_j$.  Thus $U<z_j$: properness only in
$\operatorname{Aut}(T_j)$ would not suffice.  The soluble coordinate
normal test gives $\widehat R\le D/F_j$, hence
$R\le P_{ij}\le M_{ij}$.  Its order contradicts
$q_j^{c_0n_j}>|M_{ij}|$.
\end{proof}

\subsection{Eliminating every crossing path}

Put $\mathcal E_0(x)=x$ and
$\mathcal E_{h+1}(x)=\exp(\mathcal E_h(x))$.

\begin{lemma}[Repeated crossing growth]\label{grw:towers}
Suppose every comparison in a reduced homogeneous chain has a crossing.
For every fixed $h,C$ and every sufficiently distant pair $i<j$,
\[
 n_j>\mathcal E_h(a_i^C).
\]
The required distance is a constant depending only on $h,C$ and the
constants in the catalogue.  Also $n_j>a_i^C$ for every fixed $C$ after
a fixed separation and a sufficiently large minimum rank.
\end{lemma}

\begin{proof}
Every ascending same-characteristic geometric step preserves both
natural dimension and $n\log q$ up to fixed positive lower factors,
as established in \ref{grw:input-degrees}.  The same holds for a
same-characteristic nonlinear step by its quadratic degree bound and
fixed positive field ratio.  If the last cover crosses, a preceding
same-characteristic cover leaves $n'\log q'$ at least a fixed multiple
of $n_i\log q_i$.  A preceding crossing has the same consequence
from $n'\ge\exp(c n_i\log q_i)$ and $\log q'\ge\log2$ at sufficiently
large rank.  The last crossing then gives the endpoint lower bound.
If only the first cover crosses, it gives
$n'\ge\exp(c n_i\log q_i)$, and the last cover preserves natural
dimension up to a fixed factor.  Decreasing $c$ and increasing the
fixed minimum-rank threshold, we obtain in every case
\[
 n_{i+1}\ge\exp(c n_i\log q_i).
\]
Write
$u=n_i\log q_i$.  Classical order formulas give
\[
 \log a_i\le C_1 n_i^2\log q_i\le C_2u^2.
\]
After the first step every additional comparison gives at least the
map $x\mapsto F(x)=\exp(c\log(2)x)$.  For fixed $C_3$,
\[
 \frac{F(\exp(cu))}{\exp(C_3u^2)}\longrightarrow\infty
       \qquad(u\longrightarrow\infty).
\]
If $L(u)/R(u)\longrightarrow\infty$ and $R(u)\longrightarrow\infty$,
then $F(L(u))/\exp(R(u))\longrightarrow\infty$ as well.  Induction
therefore shows that $h+2$ comparisons dominate
$\mathcal E_h(\exp(C_3u^2))$, and hence
$\mathcal E_h(a_i^C)$ when $C_3\ge CC_2$.  This also proves the last
assertion.  The target prime characteristic need not be bounded in
this argument: its only use in the iteration was $\log q\ge\log2$.
\end{proof}

\begin{proposition}[No characteristic crossings]\label{grw:no-crossing}
After passage to a reduced homogeneous chain, neither cover of any
comparison changes the defining characteristic.
\end{proposition}

\begin{proof}
Suppose otherwise, and color by the two step categories.  Every pair in
the resulting arbitrarily long chain has a crossing.  There are three
possibilities for its last step $A$.

\smallskip\noindent
\emph{The last step is geometric in the same characteristic.}
The first step $B$ crosses.  Hence $n',f'\le a_i^C$ by
\eqref{grw:cross-bounds}.  An ordinary geometric last step changes
dimension by bounded factors, except for an extension degree.  By
\eqref{grw:matrix-field-convention} that degree is
$b=\epsilon'f'/(\epsilon_jf_j)\le2f'$.  A last subfield degree can be
arbitrarily large but leaves dimension unchanged.  A twisted-tensor
last step has exponent $b=f'/f_j\le f'$ by
\eqref{grw:twisted-field-convention}.  Therefore
\[
 n_j\le (C n')^{C(1+f')}+C(1+f')n'+C
       \le\exp(a_i^{C'})
\]
for a fixed $C'$.  Lemma~\ref{grw:towers} contradicts this bound on a
sufficiently distant pair.

\smallskip\noindent
\emph{The last step is symplectic type.}
Its source $Y$ is a symplectic or orthogonal group over a prime field,
say of order $r$, and its target dimension is $r^{d}$ with
$2d+O(1)=n'$.
First suppose $B$ is in the same characteristic.  Then $r=p_i$ and
$f'=1$.  In a subfield first step the formula $1=\rho f_i b$, with
$\rho$ in a fixed finite set of positive rationals, bounds both $f_i$
and $b$; in the untwisted formula such a proper step is impossible.
An ordinary extension degree in $B$ satisfies
$b=\epsilon_i f_i\le2f_i$, because the target here has $f'=1$ and
$\epsilon'=1$; a twisted-tensor exponent satisfies $b=f_i$.
The ordinary geometric formulas give a fixed-height tower
bound on $n'$ in terms of $a_i$.  If $B$ is nonlinear, its irreducible
degree is at most $a_i^C$, giving the same conclusion.  Since
$p_i\le a_i$, the formula $n_j=p_i^{d}$ gives
$n_j\le\mathcal E_h(a_i^{C'})$ for fixed $h,C'$.  This contradicts
Lemma~\ref{grw:towers}.

It remains in this case that $B$ also crosses.  Let $b_{ij},s_{ij}$
be the inner multipliers of $B,A$.  By
\ref{grw:input-multipliers},
\[
 1\le b_{ij}\le a_i^C,
 \qquad c n_j^2\le s_{ij}\le C n_j^2,
 \qquad e_{ij}=b_{ij}s_{ij}.
\]
The bound for $b_{ij}$ is the outer-order bound when $B$ is almost
simple; in the symplectic-type case it follows from
$b_{ij}\le C(n')^2\le a_i^{C'}$.
For $i<j<k$, Lemma~\ref{grw:indices} now gives
\[
 c n_j^2\le e_{ij}
 =\frac{b_{ik}s_{ik}}{b_{jk}s_{jk}}
 \le\frac{C}{c}a_i^C.
\]
The factors involving the common upper dimension $n_k$ have canceled.
This contradicts Lemma~\ref{grw:towers} on a sufficiently distant
pair $i<j$.

\smallskip\noindent
\emph{The last step is almost simple irreducible.}
Let its source be $Y$.  If $B$ crosses, then
$n',f'\le a_i^C$ and $|P_B|\le a_i^C$: for an almost simple $B$ this
uses $|P_B|\le a_i|\operatorname{Out}(T_i)|$, and for a symplectic-type
$B$ it uses the preceding bound on its multiplier.  Equation
\eqref{grw:normalizer-input} gives
\begin{equation}\label{grw:small-M-first-cross}
 |M_{ij}|\le a_i^{C'}.
\end{equation}
It is irrelevant here whether $A$ itself is a crossing.

Suppose next that $B$ is in the same characteristic, so that $A$
crosses.  If $B$ is not a subfield step, its geometric dimension
formulas, or its irreducible degree bound when nonlinear, bound $n'$
by a fixed-height tower in $a_i$.  Its field satisfies
$\log q'\le C\log q_i$: restriction of scalars contributes the
bounded factor $\epsilon_i/\epsilon'$, and the other non-subfield
rows have their fixed field factors.  Thus $|Y|$ is bounded
by a fixed-height tower in $a_i$, and
$n_j\le |Y|^C$ contradicts Lemma~\ref{grw:towers}.

Finally suppose $B$ is subfield.  Its rank equals the lower rank and
its characteristic equals $p_i$.  Its inner projection has order at
most $a_i^C$, but $f'$ can be unbounded.  The last crossing gives
\[
 \log n_j\ge c n' f'\log p_i,
 \qquad f'\le C\log n_j.
\]
Equation~\eqref{grw:normalizer-input}, with $n'$ comparable to $n_i$,
therefore gives
\begin{equation}\label{grw:small-M-subfield}
 |M_{ij}|\le a_i^{C'}(1+\log n_j)^{C'}.
\end{equation}
The same upper bound includes \eqref{grw:small-M-first-cross}.
By Lemma~\ref{grw:towers}, after a fixed separation we have
$a_i^{C'}\le n_j^{1/16}$; and, for sufficiently large rank,
$(1+\log n_j)^{C'}\le n_j^{1/16}$.  Hence
$|M_{ij}|\le n_j^{1/8}$.  This satisfies the dimension inequality of
Lemma~\ref{grw:small-parabolic}, and
$q_j^{c_0n_j}>n_j^{1/8}\ge |M_{ij}|$.  That Lemma gives the final
contradiction.

The three possibilities exhaust the last cover in a crossing path.
\end{proof}

\subsection{A common field and the exclusion of radicals}

\begin{lemma}[Common defining field]\label{grw:common-field}
On the remaining arbitrarily long chains, the endpoints can be taken
to have one common family and one common field $q=p^f$.  The intermediate
field ratios belong to a fixed finite set of positive rationals.  In
particular, all extension and twisted-tensor degrees in these paths
are bounded integers.  Every comparison satisfies
\begin{equation}\label{grw:quadratic-growth}
 n_j\ge c n_i^2.
\end{equation}
\end{lemma}

\begin{proof}
Propositions~\ref{grw:no-geometric} and~\ref{grw:no-crossing} leave a
defining-characteristic nonlinear cover on every comparison, with the
common characteristic of the endpoints. Color by endpoint family
and cover categories.  A nonlinear cover has a bounded rational field
ratio; only the other cover can carry an exceptional prime degree.
Lemma~\ref{grw:prime-colors} applied to $f_i$ therefore gives equal
endpoint exponents.  The exceptional degree at the remaining cover
must cancel a member of a fixed finite set of rational twisting ratios,
so it is bounded.  The intermediate ratios consequently belong to a
fixed finite set.

The nonlinear cover increases natural dimension at least quadratically.
All remaining geometric steps increase it by at least a fixed positive
factor.  Thus its position as first or last cover makes no difference
to \eqref{grw:quadratic-growth}, after decreasing $c$.
\end{proof}

\begin{lemma}[No parabolic cover]\label{grw:no-parabolic-cover}
After a further homogeneous reduction, no selected shortcut contains
a parabolic or subspace cover with a large unipotent radical.
\end{lemma}

\begin{proof}
Use Lemma~\ref{grw:common-field}, and color by the position of the
parabolic cover.  The other cover is nonlinear.

If the parabolic cover is first, its active dimension is $n_i$ and its
ambient dimension is comparable to $n_i$.  Its inner quotient is at
most $q^{C n_i^2}$.  The multiplier of the last nonlinear cover is at
most $|\operatorname{Out}(Y)|$, which is absorbed by this bound, since
$n'$ is comparable to $n_i$ and $f'/f$ is bounded.  Hence
\begin{equation}\label{grw:first-radical-upper}
 e_{ij}\le q^{C_1 n_i^2}\qquad\hbox{for every }i<j.
\end{equation}
For $j<k$ the radical in its first cover also gives
$e_{jk}\ge q^{c_1 n_j}$, with the bounded field ratio absorbed in $c_1$.
By Lemma~\ref{grw:indices},
\[
 q^{c_1 n_j}\le e_{jk}\le e_{ik}\le q^{C_1 n_i^2}.
\]
Iterating \eqref{grw:quadratic-growth} through a fixed number of
intermediate chain vertices makes $n_j/n_i^2$ tend to infinity as the
minimum rank grows.  The displayed inequalities are incompatible.

If the parabolic cover is last, write $e_{ij}=b_{ij}s_{ij}$, where
$b_{ij}\le|\operatorname{Out}(T_i)|$ is the nonlinear multiplier and
$s_{ij}$ is the parabolic multiplier.  For a fixed upper endpoint $k$,
the values $s_{ik}$ have a uniformly bounded number of possibilities
by \ref{grw:input-multipliers}.  List those values in increasing order
and color $(i,k)$ by its position in that list.  This gives a uniform
finite coloring even though the lists for different $k$ may be
different.  On a homogeneous triple, $s_{ik}=s_{jk}$, and therefore
\[
 e_{ij}=\frac{e_{ik}}{e_{jk}}
       =\frac{b_{ik}}{b_{jk}}
       \le |\operatorname{Out}(T_i)|.
\]
But the parabolic radical on the pair $(i,j)$ gives
$e_{ij}\ge q^{c_2n_j}$.  The classical outer-order bound is
\[
 |\operatorname{Out}(T_i)|\le C n_i f
                         \le q^{C_2(1+\log n_i)},
\]
so even \eqref{grw:quadratic-growth} contradicts these inequalities
at sufficiently large rank.
\end{proof}

\begin{lemma}[Inner order bound]\label{grw:inner-order}
For the remaining shortcuts there is an absolute $C$ such that
\begin{equation}\label{grw:inner-small}
 |P_{ij}|\le q^{C n_i^2}.
\end{equation}
\end{lemma}

\begin{proof}
There is at least one nonlinear cover.  For a remaining geometric
cover with active label $T$ and natural dimension $n$, the other
classical factors have bounded natural dimension and are soluble.
Their orders, scalar factors, and bounded component-permutation factors
are bounded by
\[
 q^{C_0}(1+n)^{C_0}(1+f)^{C_0}.
\]
Field and twisted-tensor degrees are bounded by
Lemma~\ref{grw:common-field}; their semilinear factors satisfy the same
bound.  These are precisely the block, form, tensor, and field
structures left after Lemma~\ref{grw:no-parabolic-cover}; no unipotent
radical or extra nonabelian factor is being omitted from the estimate.

If the geometric cover is first, its intermediate natural dimension
is at most $C_1(1+n_i)^{b_1}$ for fixed constants.  Its multiplier has
the displayed bound in the lower parameters, and the last nonlinear
multiplier is at most $|\operatorname{Out}(Y)|\le Cn'f'$.
Since $|T_i|\le q^{C n_i^2}$ and $f\le q$, their product satisfies
\eqref{grw:inner-small}.

If the nonlinear cover is first, its intermediate degree is at most
$a_i^{C_2}$ and $f'/f$ is bounded.  The geometric multiplier is bounded
by $q^{C_3}(1+n')^{C_3}(1+f')^{C_3}$, which is at most
$q^{C_4 n_i^2}$.  Multiplying by $a_i$ and the first nonlinear
outer multiplier proves the same bound.  With two nonlinear covers
only the two outer multipliers occur, and the latter argument applies
without the geometric factor.
\end{proof}

\begin{proposition}[The remaining paths]\label{grw:remaining-paths}
If the reduced chains of Proposition~\ref{cls:reduction} have arbitrary
length, then they still have arbitrary length subject to all of the
following properties.
\begin{enumerate}[label=\textup{(\roman*)}]
\item The endpoints have a common classical family and field $q=p^f$,
their natural dimensions tend to infinity, and their common factor
multiplicity is $m$.  Every selected shortcut has one active factor
at each cover and no other nonabelian factor in its inner projection.

\item Every selected two-edge shortcut contains a defining-characteristic
nonlinear cover.  Its representation is a single Frobenius digit after
the twisted-tensor geometric alternative has been separated.  The
other cover is another such nonlinear cover, or a block, form, tensor,
extension-field, subfield, or twisted-tensor cover.  No crossing,
symplectic-type cover, or parabolic cover with a large unipotent radical
remains.  The characteristic-two nonsingular-line form inclusion is
retained, with its natural simple quotient convention.

\item There are fixed constants $b,C$ and a fixed finite set
$\mathcal R\subset\mathbb Q_{>0}$ such that every intermediate exponent
satisfies $f'/f\in\mathcal R$, every field-extension or twisted-tensor
degree is at most $b$, and every inactive geometric natural dimension
is at most $b$.  In particular each remaining geometric cover satisfies
$n_{\mathrm{target}}\le C(1+n_{\mathrm{source}})^b$.

\item There are fixed $c,C>0$ such that
\[
 n_j\ge c n_i^2,\qquad |P_{ij}|\le q^{C n_i^2}.
\]
For any prescribed exponent $d$, fixed-spacing thinning gives also
$n_j\ge n_i^d$ for every retained pair, after raising the minimum rank.

\item The coordinate normal tests of
Lemma~\ref{grw:coordinate-normal-tests} remain available for every tested proper
coordinate overgroup.  After one more fixed-spacing thinning,
$M_{ij}$ normalizes no proper parabolic of $T_j$.
\end{enumerate}
\end{proposition}

\begin{proof}
Parts (i)--(iv) collect Propositions~\ref{grw:no-geometric} and
\ref{grw:no-crossing} and Lemmas~\ref{grw:common-field},
\ref{grw:no-parabolic-cover}, and \ref{grw:inner-order}.
For the spacing assertion put $x_i=c n_i$.  Then
\eqref{grw:quadratic-growth} says $x_j\ge x_i^2$ on every consecutive
comparison.  After $s$ comparisons,
$n_j\ge c^{-1}(c n_i)^{2^s}$.  Choose $2^s>d$ and then increase the
minimum rank to obtain $n_j\ge n_i^d$.

The coordinate normal test is inherited from the original tested
intervals and is not altered by deleting chain vertices.  Every proper
parabolic of a large-rank classical group has unipotent radical of
order at least $q^{c_3 n_j}$ for an absolute $c_3>0$.
One can see the uniform bound by choosing a maximal parabolic
containing it: the omitted simple root has at least the relative rank
many positive-root coordinates outside the Levi subsystem.  Their
product contributes at least $q^{c_3 n_j}$; unitary root orbits can
only increase the corresponding root-group orders.  The radical of
the smaller parabolic contains that of the maximal one.
Fixed-spacing thinning makes $c_3n_j>Cn_i^2$ uniformly.
If $M_{ij}$ normalized such a parabolic $J$ with radical $R$, use
\eqref{grw:normal-enlargement}.  Since
$P_{ij}R\le N_{T_j}(J)<T_j$, its full inverse image is proper in
$z_j$.  The soluble product $\widehat R\trianglelefteq\overline U$
then satisfies \eqref{grw:normal-containment}.  In particular
$R\le P_{ij}$,
contrary to $|R|\ge q^{c_3 n_j}>q^{C n_i^2}\ge|P_{ij}|$.
This proves (v).
\end{proof}

The Proposition makes no assertion that a unique nonabelian composition
factor alone produces a quasisimple subgroup.  The actual components
and their compatible rational maps will be obtained from the retained
geometric structures in the next section.  In particular, the
characteristic-two form case is still among those structures and
requires its own form-preservation argument.
